% Standalone problem document, generated from the adjacent README.md.
% Compile with XeLaTeX. Mathematical content is maintained in Markdown.
\documentclass[11pt,a4paper]{article}
\usepackage[fontsize=10.5pt]{scrextend}
\usepackage[margin=22mm,headheight=15pt,headsep=9mm,footskip=11mm]{geometry}
\usepackage{fontspec}
\setmainfont{texgyrepagella}[Extension=.otf,UprightFont=*-regular,BoldFont=*-bold,ItalicFont=*-italic,BoldItalicFont=*-bolditalic]
\setsansfont{texgyreheros}[Extension=.otf,UprightFont=*-regular,BoldFont=*-bold,ItalicFont=*-italic,BoldItalicFont=*-bolditalic]
\setmonofont{texgyrecursor}[Extension=.otf,UprightFont=*-regular,BoldFont=*-bold,ItalicFont=*-italic,BoldItalicFont=*-bolditalic,Scale=MatchLowercase]
\usepackage{amsmath,amssymb}
\usepackage{unicode-math}
\setmathfont{texgyrepagella-math.otf}
\usepackage{xcolor,microtype,enumitem,needspace,fancyhdr}
\usepackage{bookmark}
\definecolor{ink}{HTML}{17344B}
\definecolor{accent}{HTML}{197D80}
\definecolor{muted}{HTML}{596773}
\hypersetup{colorlinks=true,linkcolor=accent,urlcolor=accent,pdftitle={MI-16 - The
maximum permanent on a positive-semidefinite unitary
orbit},pdfauthor={Open Problems in Numerical Linear Algebra}}
\urlstyle{same}
\setlength{\parindent}{0pt}
\setlength{\parskip}{5pt plus 1pt minus 1pt}
\linespread{1.04}
\setlength{\emergencystretch}{3em}
\widowpenalty=10000
\clubpenalty=10000
\displaywidowpenalty=10000
\allowdisplaybreaks[1]
\setlist{nosep,leftmargin=1.5em}
\providecommand{\tightlist}{\setlength{\itemsep}{0pt}\setlength{\parskip}{0pt}}
\providecommand{\pandocbounded}[1]{#1}
\pagestyle{fancy}
\fancyhf{}
\fancyhead[L]{\sffamily\footnotesize\color{muted}OPEN PROBLEMS IN NUMERICAL LINEAR ALGEBRA}
\fancyhead[R]{\sffamily\bfseries\footnotesize\color{accent}MI-16}
\fancyfoot[L]{\parbox[b]{0.9\textwidth}{\sffamily\fontsize{8}{10}\selectfont\color{muted}Editorial ratings; a bounded search cannot certify that a problem remains unsolved.\\Literature check: 2026-09-12}}
\fancyfoot[R]{\sffamily\footnotesize\color{muted}\thepage}
\renewcommand{\headrulewidth}{0.3pt}
\renewcommand{\headrule}{\hbox to\headwidth{\color{accent}\leaders\hrule height \headrulewidth\hfill}}
\makeatletter
\renewcommand{\section}{\Needspace{5\baselineskip}\@startsection{section}{1}{0pt}{12pt plus 2pt minus 2pt}{4pt}{\sffamily\bfseries\large\color{ink}}}
\renewcommand{\subsection}{\Needspace{4\baselineskip}\@startsection{subsection}{2}{0pt}{9pt plus 1pt minus 1pt}{3pt}{\sffamily\bfseries\normalsize\color{ink}}}
\makeatother
\setcounter{secnumdepth}{0}
\begin{document}
{\sffamily\small\color{muted}Matrix inequalities and norms\par}
\vspace{3pt}
{\raggedright\hyphenpenalty=10000\exhyphenpenalty=10000\sffamily\bfseries\fontsize{21}{24}\selectfont\color{ink}The
maximum permanent on a positive-semidefinite unitary orbit\par}
\vspace{7pt}
{\sffamily\small\textbf{Difficulty:} extreme\quad\textbf{Importance:} interesting
to the community\par}
{\sffamily\small\bfseries\color{accent}Status: Solved\par}
\vspace{4pt}
{\color{accent}\hrule height 0.8pt}
\vspace{6pt}
\textbf{Provenance:} explicit exact extremal problem; no conjectured
optimizer is supplied

\textbf{Historical rating rationale:} An exact formula for arbitrary
spectra would resolve a longstanding permanent optimization barrier and
inform prescribed-spectrum matrix optimization.

\subsection{Resolution --- exact algebraic determination, 12 September
2026}\label{resolution-exact-algebraic-determination-12-september-2026}

\textbf{Solved by Sidney Holden}, Center for Computational Biology,
Flatiron Institute, Simons Foundation.
\href{https://github.com/ajt60gaibb/OpenProblemsInNLA/blob/main/references/holden-mi16-2026-09-12/README.md}{Submission
and verified affiliation}.
\href{https://github.com/ajt60gaibb/OpenProblemsInNLA/blob/main/matrix-inequalities-and-norms/MI-16/solution.pdf}{Proof
PDF} ·
\href{https://github.com/ajt60gaibb/OpenProblemsInNLA/blob/main/matrix-inequalities-and-norms/MI-16/solution.tex}{LaTeX
source} ·
\href{https://github.com/ajt60gaibb/OpenProblemsInNLA/blob/main/references/holden-mi16-2026-09-12/independent-review.md}{Independent
AI-agent audit: PASS}.

Theorem 2.1 and Sections 3--5 give an exact finite algebraic
prescription for every nonnegative real spectrum in every order. A
specialized critical-value elimination polynomial gives finitely many
candidates. Explicit Schur-character sums compute a spectral moment at a
proved finite order, which selects the maximum even when some real
critical values are not Hermitian-feasible. The prescription has no
remaining matrix optimization or unevaluated limit, and meets the
original exact-value criterion below. For arbitrary real inputs it uses
exact ordered-field arithmetic and comparisons; the reference programs
accept rational inputs.

\textbf{Form of the answer:} this is not a compact general structural
formula, an efficient arbitrary-order solver, or a classification of all
general-spectrum optimizers. Those stronger goals are not claimed.
Theorem 8.1 separately gives a two-candidate structural formula for one
exceptional eigenvalue; Sections 9--10 prove the transition threshold
and every equality case. The earlier rank-one expansion and
equal-support reduction from
\href{https://github.com/ajt60gaibb/OpenProblemsInNLA/pull/186}{Holden's
partial submission, PR \#186}, are retained and credited in the dossier.
This submission adds the all-spectrum determination and sharper
exceptional-spectrum results.

A separate Codex agent independently audited the complete argument and
its correspondence with this unchanged target. This is informal AI-agent
review, not external human peer review or formal verification. No Lean
verification or historical-priority claim is made. Whether a structural
formula is preferred for the historical problem is exposed for
maintainer review; it is not silently added to this entry's completion
criterion.

\subsection{Problem statement}\label{problem-statement}

For every integer \(n\ge1\) and every real list
\(\lambda=(\lambda_1,\ldots,\lambda_n)\) with \(\lambda_i\ge0\),
determine the exact value of

\[M_n(\lambda)=\max_{U\in\mathcal U_n}
\mathop{\mathrm{per}}\nolimits\!\left(U^*\mathop{\mathrm{diag}}\nolimits(\lambda_1,\ldots,\lambda_n)U\right),
\qquad \mathcal U_n=\{U\in\mathbb C^{n\times n}:U^*U=I_n\},\]

where
\(\mathop{\mathrm{per}}\nolimits(C)=\sum_{\sigma\in S_n}\prod_{i=1}^n c_{i,\sigma(i)}\)
and \(S_n\) is the permutation group. The permanent of these Hermitian
positive-semidefinite matrices is real. Compactness ensures the maximum
exists. A resolution should express its value from the prescribed
eigenvalues for arbitrary order, rather than only bound the objective or
restate its optimization definition.

\subsection{Relevance}\label{relevance}

The eigenvalues fix an entire unitary similarity class while the
permanent varies within it. Determining its extreme value is a concrete
optimization problem over matrices with a prescribed spectrum.

\subsection{References}\label{references}

\begin{enumerate}
\def\labelenumi{\arabic{enumi}.}
\tightlist
\item
  F. Zhang, \emph{An update on a few permanent conjectures}, Special
  Matrices 4 (2016), 305--316, passage headed ``Marcus--Minc
  max-per-\(U\) problem 1965.''
  \href{https://arxiv.org/html/1608.02844v1}{Primary text};
  \href{https://doi.org/10.1515/spma-2016-0030}{journal}.
\item
  J. H. Drew and C. R. Johnson, \emph{Counterexample to a conjecture of
  Mehta regarding permanental maximization}, Linear and Multilinear
  Algebra 25(3) (1989), 253--254, the counterexample.
  \href{https://doi.org/10.1080/03081088908817948}{DOI}.
\end{enumerate}

\subsection{Status check --- 2026-09-10}\label{status-check-2026-09-10}

Zhang's latest arXiv record is v1 and explicitly lists the general
maximization problem. Its statement and discussion of the equal-diagonal
counterexample were inspected. The original 1989 full text was not
retrieved; that counterexample's role was checked through Zhang's
explicit account and bibliography. Searches used
\textit{maximum permanent unitary 2026},
\textit{maximum permanent eigenvalues 2025 2026}, and
\textit{Marcus-Minc max-per-U problem}. No general spectral formula
was located. The claim that an equal-diagonal representative always
maximizes the permanent is false and is not imposed. The other catalog
permanent entries concern inequalities between distinct matrix functions
or products, not this exact orbit maximum.

\textbf{Audit update (2026-09-10):} Rechecked Zhang's Marcus--Minc
max-per-\(U\) statement and searched for later exact spectral formulas.
Equal-diagonal optimality remains an invalid shortcut; no all-spectrum
determination was located. This is a bounded literature check, not a
proof that no solution exists.
\end{document}
